%% kapa: exam=1 round=29 year=2018 period=1 form=A subject=civil_law subject_ko=민법 seq=15 number=15
%% source: Q-Net 감정평가사 기출문제 (cst00309) artlSeq=5202808 file=extracted/1008/5202808/2018년 제29회 감정평가사 1차시험 기출문제/2018년 제29회 감정평가사 1차시험 기출문제/1교시 A형.pdf
%% conversion: kordoc 4.24.51 -> tools/gen1.py (auto); review=proofread-sonnet; flags=none
\begin{question}{15}
법률행위의 취소에 관한 설명으로 옳지 않은 것은? (다툼이 있으면 판례에 따름)
\begin{choices}
  \item 법률행위를 취소한 후라도 무효행위 추인의 요건을 충족할 경우, 무효행위의 추인은 가능하다.
  \item 제한능력자가 맺은 계약은 추인이 있을 때까지 상대방이 그 의사표시를 취소할 수 있다.
  \item 제한능력을 이유로 법률행위가 취소된 경우, 제한능력자는 그 행위로 인하여 받은 이익이 현존하는 한도에서 상환할 책임이 있다.
  \item 법률행위의 취소를 전제로 한 소송상의 이행청구에는 취소의 의사표시가 포함되어 있다고 볼 수 있다.
  \item 취소권은 추인할 수 있는 날로부터 3년내에 법률행위를 한 날로부터 10년내에 행사하여야 한다.
\end{choices}
\end{question}
\answer{2}
